\documentclass[11pt]{article}
\usepackage{amsmath,amssymb,amsthm}
\usepackage[margin=1in]{geometry}

\newcommand{\R}{\mathbb{R}}
\newcommand{\Vol}{\operatorname{Vol}}
\newcommand{\conv}{\operatorname{conv}}

\theoremstyle{definition}
\newtheorem{definition}{Definition}
\newtheorem{proposition}{Proposition}

\title{Primary and Secondary Voronoi Cell Types:\\
A Classification via the Projected Secondary Polytope}
\author{}
\date{}

\begin{document}
\maketitle

\section{Background}

Consider a crystal structure with atoms at positions $\{x_i\}$.
The Voronoi cell of a central atom partitions into faces, edges, and
vertices.  A vertex of the Voronoi cell is the circumcenter of a
Delaunay tetrahedron; it is equidistant from the four atoms of that
tetrahedron.  A vertex is \emph{degenerate} (or \emph{unstable}) when
$k > 4$ atoms are equidistant from it.  FCC and HCP crystals have
degenerate vertices shared by $k = 6$ Voronoi cells (the equidistant
atoms form an octahedron), while simple cubic (SC) has degenerate
vertices shared by $k = 8$ cells (the atoms form a cube).

Under infinitesimal perturbation of the atom positions, each degenerate
vertex \emph{resolves}: the single degenerate point breaks into several
generic Voronoi vertices, each equidistant from exactly four atoms.
Different perturbation directions yield different resolutions, and hence
different Voronoi cell topologies.  The family $\mathcal{F}_\lambda$ of
all topological types obtainable by infinitesimal perturbation of a
crystal $\lambda$ was introduced in~\cite{lazar2015}.

The distinction between \emph{primary} and \emph{secondary} types was
established empirically in~\cite{lazar2015}: primary types have positive
ideal frequency $f_\lambda(\tau) > 0$, meaning they are observed with
finite probability as the perturbation magnitude $\varepsilon \to 0$,
while secondary types have $f_\lambda(\tau) = 0$.  For FCC, 44 primary
types and 6{,}250 secondary types were identified, and these groupings
were validated by molecular dynamics simulation.

In~\cite{lazar2015}, the primary/secondary classification for FCC and
HCP was performed by direct geometric analysis of the octahedral vertex:
each unstable vertex resolves into either an edge (two ways) or a square
face, and all other resolutions were classified as secondary.  Here we
provide a general, computable criterion that extends this classification
to arbitrary point configurations, using the theory of secondary
polytopes.


\section{Resolutions and Regular Triangulations}

Let $p_1, \ldots, p_k \in \R^3$ be the $k$ atoms equidistant from a
degenerate Voronoi vertex $v$, with $p_1$ being the central atom.  Under
a small perturbation $\delta_i \in \R^3$ applied to each atom, the
Delaunay triangulation of the perturbed points $q_i = p_i + \delta_i$
determines a \emph{resolution} of the degenerate vertex.

The \emph{star} of the central atom is the set of Delaunay tetrahedra
containing $p_1$.  Two perturbations that produce the same star yield the
same local Voronoi cell topology at that vertex.  We call each distinct
star a \emph{resolution type} or \emph{star type}.

The Delaunay triangulation of the perturbed points is a \emph{regular
triangulation} of the original point configuration $A = \{p_1, \ldots,
p_k\}$.  By the theory of Gel'fand, Kapranov, and
Zelevinsky~\cite{gkz1994}, the set of all regular triangulations of $A$
can be parameterized by \emph{height functions} $\omega =
(\omega_1, \ldots, \omega_k) \in \R^k$, where the triangulation induced
by $\omega$ is the projection of the lower convex hull of
$\{(p_i, \omega_i)\} \subset \R^4$.  Two height vectors $\omega$ and
$\omega'$ that differ by an affine function of the coordinates (i.e.,
$\omega_i - \omega_i' = a \cdot p_i + b$ for some $a \in \R^3$,
$b \in \R$) induce the same triangulation.  The effective parameter space
is therefore the quotient $\R^k / \text{Aff}(A)$, which has dimension
$k - 4$ when $A$ affinely spans~$\R^3$.


\section{The GKZ Secondary Polytope}

The \emph{GKZ secondary polytope} $\Sigma(A) \subset \R^k$ is a convex
polytope whose vertices are in bijection with the regular triangulations
of~$A$.  For a triangulation $T$, its \emph{GKZ vector}
$\varphi(T) \in \R^k$ has components
\begin{equation}
  \varphi_i(T) = \sum_{\substack{\sigma \in T \\ p_i \in \sigma}}
  \Vol(\sigma),
  \label{eq:gkz}
\end{equation}
where the sum is over all maximal simplices $\sigma$ of $T$ that
contain~$p_i$, and $\Vol(\sigma)$ is the Euclidean volume of the
simplex.

The \emph{normal fan} of $\Sigma(A)$ is the \emph{secondary fan}, which
partitions $\R^k$ into polyhedral cones.  The normal cone of each vertex
$\varphi(T)$ is the \emph{secondary cone}
\[
  \sigma(T) = \bigl\{\omega \in \R^k : \text{$\omega$ induces
  triangulation } T\bigr\},
\]
which is a full-dimensional cone in~$\R^k$.

After projecting to the quotient $\R^{k-4} \cong \R^k / \text{Aff}(A)$,
we obtain the \emph{projected secondary polytope}
$\bar\Sigma(A) \subset \R^{k-4}$.  Since the affine span of $\Sigma(A)$
is a translate of $\text{Aff}(A)^\perp$, this projection is a
combinatorial isomorphism: every vertex of $\Sigma(A)$ maps to a vertex
of~$\bar\Sigma(A)$.


\section{Classification Criterion}

We now state the classification criterion.

\begin{definition}
A triangulation of~$A$ is \emph{proper} if every maximal simplex has
positive volume.  A star type (resolution)~$S$ is \emph{primary} if at
least one proper regular triangulation of~$A$ produces star~$S$;
otherwise $S$ is \emph{secondary}.
\end{definition}

\begin{proposition}
The GKZ secondary polytope $\Sigma(A)$ has one vertex for each regular
triangulation of~$A$.  The affine span of $\Sigma(A)$ is a translate of
$\operatorname{Aff}(A)^\perp$, so the projection
$\Sigma(A) \to \bar\Sigma(A) \subset \R^{k-4}$ is a combinatorial
isomorphism: every vertex of $\Sigma(A)$ maps to a vertex of
$\bar\Sigma(A)$.  Consequently, every proper regular triangulation has a
full-dimensional normal cone in the effective height space $\R^{k-4}$,
meaning it is induced by an open set of effective height functions and
arises under a positive-measure set of perturbation directions as
$\varepsilon \to 0$.
\end{proposition}

Primary star types therefore have positive solid angle in the effective
height space.  A global Voronoi cell type is primary if the resolution
at every degenerate vertex is a primary star type.

\subsection{Degenerate triangulations and secondary types}

The flip-graph enumeration of regular triangulations can produce
configurations that include \emph{zero-volume simplices}: tetrahedra
whose four vertices are co-planar.  These arise when a lower-dimensional
face of the convex hull $\conv(A)$ admits more than one triangulation,
and the enumeration explores subdivisions of such faces.  We call a
triangulation containing any zero-volume simplex \emph{degenerate}.

For the octahedral configuration ($k = 6$), the flip-graph enumeration
produces 9~triangulations, of which only 3~are proper.  For the cubic
configuration ($k = 8$), the enumeration produces
11{,}648~triangulations, of which only 74~are proper.

Degenerate triangulations are problematic for the primary classification
in two ways:
\begin{enumerate}
\item \textbf{GKZ vector inflation.}  A zero-volume simplex contributes
  nothing to the GKZ vector~\eqref{eq:gkz}, so many degenerate
  triangulations share their GKZ vector with a proper triangulation that
  differs only in the subdivision of co-planar faces.  For the cube,
  11{,}648~triangulations produce only 152~distinct GKZ vectors, while
  the 74~proper triangulations have 74~distinct vectors.
\item \textbf{Phantom star types.}  A zero-volume simplex that contains
  the central atom~$p_1$ changes the star without changing the GKZ
  vector.  This creates star types that are not realizable by any proper
  triangulation---they require co-planar atoms to remain exactly
  co-planar under perturbation, a measure-zero condition.  For the cube,
  the 544~total star types include 512~phantom types that are only
  realized by degenerate triangulations.
\end{enumerate}

The resolution is to filter to proper triangulations before performing
the GKZ analysis.  By the GKZ bijection theorem, proper regular
triangulations have unique GKZ vectors, and the projected secondary
polytope faithfully represents their combinatorial structure.

\begin{paragraph}{Regularity of proper triangulations.}
The algorithm enumerates regular triangulations via flip-graph BFS,
which is complete over the regular triangulation graph by the GKZ
connectivity theorem.  A natural question is whether there exist proper
triangulations that are \emph{not} regular and hence missed by this
enumeration.  For $k \leq d + 3 = 6$ points in $\R^3$, all
triangulations are regular~\cite{dlrs2010}.  For the cube ($k = 8$), we
verify computationally that all 74~proper triangulations appear in the
flip-graph enumeration.  In general, non-regular triangulations of
convex point sets exist for sufficiently large~$k$, but this does not
affect correctness: only regular triangulations can arise from
height-function perturbations, so any non-regular proper triangulation
(if one existed) would not be physically realizable and is correctly
excluded.
\end{paragraph}


\section{Computational Procedure}

The classification can be computed as follows for a degenerate vertex
with $k$ equidistant atoms $A = \{p_1, \ldots, p_k\}$:

\begin{enumerate}
\item \textbf{Enumerate regular triangulations.}
  Use a flip-graph BFS to find all regular triangulations of~$A$.
  This enumeration is complete by the GKZ connectivity theorem.

\item \textbf{Filter to proper triangulations.}
  Discard any triangulation containing a zero-volume simplex (four
  co-planar vertices).  These degenerate triangulations arise from
  subdivisions of lower-dimensional faces and represent measure-zero
  perturbation conditions.  Let $T_1, \ldots, T_M$ denote the
  remaining proper triangulations.

\item \textbf{Compute GKZ vectors.}
  For each proper triangulation $T_j$, compute $\varphi(T_j) \in \R^k$
  using equation~\eqref{eq:gkz}.  By the GKZ bijection theorem, these
  $M$~vectors are distinct.

\item \textbf{Project to $\R^{k-4}$.}
  Let $V \subset \R^k$ be the affine span of $A$ augmented by the
  constant vector (dimension~4).  Project each GKZ vector onto
  $V^\perp$ to obtain points in $\R^{k-4}$.

\item \textbf{Identify hull vertices.}
  Compute the convex hull of the $M$ projected points.  For all
  configurations tested, every proper triangulation projects to a vertex
  of this hull (consistent with the combinatorial isomorphism between
  $\Sigma(A)$ and $\bar\Sigma(A)$).

\item \textbf{Extract primary star types.}
  For each hull vertex, extract the star of the central atom from the
  corresponding triangulation.  The union of these stars gives the set
  of primary star types.
\end{enumerate}


\section{Validation and Results}

\subsection{FCC octahedral vertex ($k = 6$)}

The six atoms equidistant from an octahedral vertex in FCC form a
regular octahedron.  The flip-graph enumeration produces 9~triangulations,
of which 3~are proper (the remaining 6~contain zero-volume tetrahedra
from co-planar atom triples).  The 3~proper triangulations yield 3~distinct
star types, all of which are primary:
\begin{itemize}
  \item Two \emph{edge} resolutions: the degenerate vertex becomes an
    edge of the Voronoi cell (the far atom does not become a neighbor).
  \item One \emph{square} resolution: a new quadrilateral face forms
    between the central atom and the far atom.
\end{itemize}
The 6~degenerate triangulations produce 4~additional star types (in
which the far atom becomes a neighbor via a triangular face); these are
secondary.  This recovers exactly the classification
from~\cite{lazar2015}, where the 44~primary FCC types arise from
$3^6 = 729$ combinations over the six unstable vertices.

\subsection{SC cubic vertex ($k = 8$)}

The eight atoms equidistant from a vertex of the simple cubic Voronoi
cell form a cube.  The flip-graph enumeration produces
11{,}648~triangulations, of which only \textbf{74~are proper} (the
remaining 11{,}574 contain zero-volume tetrahedra from co-planar atom
quadruples on the cube faces).  The 74~proper triangulations have
74~distinct GKZ vectors, all of which are vertices of the projected
secondary polytope in $\R^{8-4} = \R^4$.  These 74~proper
triangulations produce \textbf{32~distinct star types}, all primary.
The total number of star types (including degenerate triangulations)
is~544, so 512~are secondary.

This yields a substantial reduction in the combinatorial enumeration
space.  With 8~degenerate vertices per SC Voronoi cell:
\begin{center}
\begin{tabular}{lcc}
  \hline
  & Resolutions/vertex & Total combinations \\
  \hline
  All types & 544 & $544^8 \approx 7.7 \times 10^{21}$ \\
  Primary only & 32 & $32^8 \approx 1.1 \times 10^{12}$ \\
  \hline
\end{tabular}
\end{center}
The reduction factor is approximately $10^{10}$, far more dramatic
than the FCC case.  The $O_h$~point group symmetry ($|G| = 48$) and
validity pruning provide further reduction.


\begin{thebibliography}{9}

\bibitem{lazar2015}
E.~A.~Lazar, J.~Han, and D.~J.~Srolovitz,
``Topological framework for local structure analysis in condensed
matter,'' \emph{Proc.\ Natl.\ Acad.\ Sci.\ USA}, vol.~112, no.~43,
pp.~E5769--E5776, 2015.

\bibitem{gkz1994}
I.~M.~Gel'fand, M.~M.~Kapranov, and A.~V.~Zelevinsky,
\emph{Discriminants, Resultants, and Multidimensional Determinants}.
Birkh\"auser, 1994.

\bibitem{dlrs2010}
J.~A.~De~Loera, J.~Rambau, and F.~Santos,
\emph{Triangulations: Structures for Algorithms and Applications}.
Springer, 2010.

\end{thebibliography}

\end{document}
